% %%%%%%%%%%%%%%%%%%%%%%%%%%%%%%%%%%%%%%%%%%%%%%%%%%%%%%%%%%%%%%%%%%%%%%%%%%%%%%%%%
%
%	cap#.tex 
%  -----------
%
%	Neste arquivo você escreve o capítulo. Deve sempre iniciar com o ambiente
%	\chapter{Nome do Capitulo}
%
% %%%%%%%%%%%%%%%%%%%%%%%%%%%%%%%%%%%%%%%%%%%%%%%%%%%%%%%%%%%%%%%%%%%%%%%%%%%%%%%%%
%
\chapter{Conclusões}

Este trabalho teve como objetivo principal investigar a influência de galerias transversais nas deformações em túneis gêmeos profundos, através de uma análise computacional considerando diferentes modelos de comportamento para o revestimento e para o maciço. A partir dessa abordagem, foi possível investigar a distribuição de tensões, o comportamento das deformações e a interação entre os túneis em cenários com e sem a presença de uma galeria transversal, além da influência do parâmetro friccional.

Sob a perspectiva constitutiva, a parcela irreversível da deformação do maciço rochoso é representada por meio do modelo acoplado elastoplástico-viscoplástico. Essa abordagem se mostra adequada especialmente para descrever tanto as deformações imediatas quanto aquelas que ocorrem de forma retardada em rochas argilosas profundas.



A variação de $\rho(\theta)$ ao longo do perímetro da seção dos túneis sem galeria revela uma anisotropia das deformações, indicando que o comportamento da parede do túnel não é uniforme, como ocorre em túneis isolados. Tal comportamento é mais acentuado no longo prazo, principalmente para baixos ângulos de atrito.



